\documentclass{situationist-preprint}

\title{The Symmetry Objection Answered}
\msid{EPI-04-b}
\msfamily{EPI-04}
\msdiscipline{Epistemology}
\mstype{Consequence}
\mskeywords{grue; bleen; projectibility; naturalness; language-relativity; Goodman}

\begin{document}
\maketitle

\begin{abstract}
Goodman's symmetry objection holds that whether a predicate is qualitative depends on the language: in a language with \enquote{grue} and \enquote{bleen} as primitives, \enquote{green} is the positional predicate. We derive from the thesis of EPI-04-a, that projectibility tracks naturalness measured by the shortest construction from the sparse universals, that the objection fails. The measure is fixed by the universals, not by any language; the construction of grue requires a particular time as a parameter whatever language describes it; and the grue speaker's definition of green, though short in her primitives, is long in the measure that matters.
\end{abstract}

\section{The objection}

An object is \emph{grue} if it is green and examined before $t$, or blue and not so examined; it is \emph{bleen} if it is blue and examined before $t$, or green and not so examined. \citet[ch.~3]{Goodman1955} observed that \enquote{green} can be defined from \enquote{grue} and \enquote{bleen} exactly as \enquote{grue} is defined from \enquote{green} and \enquote{blue}: an object is green if it is grue and examined before $t$, or bleen and not so examined. If projectibility were a matter of a predicate's being purely qualitative rather than positional, it would depend on which predicates a language takes as primitive, and the grue speaker would have as good a claim as we do. EPI-04-a gave the answer in brief (its Section~3.3); we give it in full.

\section{The answer}

EPI-04-a measures the naturalness of a property by the complexity of its shortest construction from the sparse universals, counting particular objects used as parameters (its Definition~2.1), and holds that projectibility tracks naturalness (Principle~P1). The sparse universals are intrinsic properties and relations: they do not distinguish intrinsic duplicates \citep{Lewis1983}.

\begin{proposition}\label{p:indep}
The naturalness of a property does not depend on the language in which it is expressed.
\end{proposition}

\begin{proof}
A construction is a way of building the property from the universals, and its complexity is counted in universals, operations and parameters. A language enters only in naming the property; two predicates of different languages that denote the same property have the same shortest construction.
\end{proof}

\begin{proposition}\label{p:param}
Every construction of grue uses a particular time, or another particular, as a parameter. Some construction of green uses none.
\end{proposition}

\begin{proof}
By Proposition~3.1 of EPI-04-a, grue distinguishes intrinsic duplicates: of two intrinsically duplicate green emeralds, one examined before $t$ and one not, the first is grue and the second is not. A construction from intrinsic universals using no particulars yields only properties shared by intrinsic duplicates. Hence every construction of grue uses a particular. Green does not distinguish intrinsic duplicates, and is constructed from the universals of color, or is one of them, without parameters.
\end{proof}

The grue speaker's definition of green has the same length as ours of grue when counted in her primitives. Counted in universals it is longer than her definition of grue, since each of \enquote{grue} and \enquote{bleen} expands into a construction containing the parameter $t$, and the definition of green from them uses both. Proposition~\ref{p:indep} shows that the second count, not the first, is the one Definition~2.1 of EPI-04-a uses. Corollary~3.2 of EPI-04-a, that green is more natural and more projectible over emeralds than grue, therefore holds in every language, including the grue speaker's.

\section{Objection and reply}

\begin{objection}
Which properties are sparse universals is itself fixed relative to our classifications, so the language-independence is only apparent.
\end{objection}

\begin{reply}
The sparse universals are those whose sharing makes for objective resemblance and in terms of which the laws are stated (Principles~P3 and~P4 of EPI-04-a). Which properties these are is a question for physics, not for our predicates; the grue speaker, if she is right about the laws, must state them with the same universals we do. Goodman's symmetry holds between predicates; it does not hold between the properties they denote.
\end{reply}

\begin{thebibliography}{9}\small
\bibitem[Goodman(1955)]{Goodman1955} Goodman, N. (1955). \emph{Fact, Fiction, and Forecast}. Harvard University Press.
\bibitem[Lewis(1983)]{Lewis1983} Lewis, D. (1983). New work for a theory of universals. \emph{Australasian Journal of Philosophy}, 61(4), 343--377.
\end{thebibliography}

\end{document}
